% Modernised reading — fonds Alexandre Grothendieck, shelfmark 16, pages 1-55.
%
% This is an INTERPRETATION, not a transcription. It re-reads the pages in
% today's notation and vocabulary, states the hypotheses the manuscript leaves
% implicit, and drops the critical apparatus so that the mathematics reads
% straight through. Everything the brackets and underlines carried moves into a
% footnote. It is held to being mathematically correct as it stands; where the
% manuscript is loose or mistaken, the true statement is what appears here and
% the footnote says what the page has. In any dispute of substance the
% transcription governs: cite batch-01/02/03.fr.tex.
%
% Pass: Opus 5 (claude-opus-5), 2026-09-19 — first pass, unchecked against the
% pages by a human.
%
% Scope: one document for the shelfmark, its three batches read as one
% argument. Thirty-nine of the fifty-five leaves carry the folder's subject;
% the rest is unrelated paper he wrote on the back of (an EGA IV galley, a
% Bourbaki « Tribu », typescripts on the Brauer group, on group schemes and on
% Banach categories) and is not mentioned. Pages 18 and 20 are two leaves of
% another text of his, on the Picard scheme and l-adic cohomology, and are
% given one paragraph, marked as foreign.
%
% Spine. The folder asks one question three times over, at three altitudes.
% (1) pp. 2-17, « Sur certains anneaux gradués »: PURE ALGEBRA. Given a graded
% ring E and an element L of degree 2, when is the grading itself — the family
% of projectors pi_i — already determined by (E, L), and when do L's inverse
% Lambda and the projectors of the primitive decomposition lie in E? The
% theorem of his no. 6 answers it; his no. 7 builds the universal ring
% Phi_0 = Z[L_0, Lambda_0] on an exterior algebra and makes the conditions a
% graded ring homomorphism Phi_0 -> E; a Question answered « Non » shows L
% alone never suffices. (2) pp. 19-28: EXAMPLES AND COUNTEREXAMPLES. The
% attempt to write Lambda as a polynomial in L (p. 19); a six-dimensional
% algebra in which L^n =/= 0 but the Lefschetz property fails (p. 21); and the
% « algèbres de Hodge d'ordre 2 » (pp. 24-28), an axiomatisation of the
% cohomology ring of a surface as a graded Poincaré-duality algebra with a
% Lefschetz element. (3) pp. 32-54: THE STANDARD CONJECTURES. The formulary of
% L and Lambda along a smooth hyperplane section (pp. 32-39), the equivalence
% C(X) <=> [C(Y) and A^0(X x Y)] (p. 41), two letters to Coates of January 1967
% and January 1966 (pp. 43-49), and the Lemma (pp. 50-54) that carries the
% theorem of no. 6 into the ring of algebraic correspondences.
%
% What only a folder-wide reading can say: run (3)'s p^X_i IS run (1)'s
% pi_{i,0} — the same projectors under two names, so that the formulary's
% identities (4) and (5) are the theorem of no. 6 read in cohomology; and the
% Lemma of p. 50 is the existence half of no. 5 and no. 6 with « stable par
% composantes homogènes » doing the work the over-ring F did on p. 7. Page 21's
% counterexample is the answer to the chain of necessary conditions p. 17
% breaks off on. None of these three links is stated on any single leaf.
%
% NOTATION fixed for the folder. Three systems collide and are separated here:
%   pi_i     (pp. 2-17) projection of M onto M^i;
%   pi_{i,alpha} (pp. 8-17) projection onto L^alpha P^{i-2alpha};
%   p^X_i    (pp. 32-39) projection onto the PRIMITIVE part P^i(X) for i <= n,
%            and onto the top L^{n-i'}P^{i'} (i = 2n - i') for i > n — i.e.
%            exactly pi_{i,0} of the first run. Written p^{X}_{i} throughout.
%   p_i      (pp. 50-54) projection onto H^i(X). Written p_i, and the reading
%            says at each use which of the three is meant.
% L is the Lefschetz operator (degree 2), Lambda its PSEUDO-INVERSE of degree
% -2 — not the Kähler adjoint; see the footnote at no. 6. E for his script E,
% Phi_0 = Z[L_0, Lambda_0], M for the graded module, P^i for primitive parts.
% phi : Y -> X is the inclusion of the hyperplane section (the typescript draws
% a slashed O), with phi^* and phi_*; xi_X, xi_Y the polarisation classes;
% chi the index of the cohomology theory. The phi, psi, Q, chi of the Hodge
% algebras (pp. 24-28) are forms on V and W and have nothing to do with the
% inclusion phi — the collision is his and is flagged in place.
%
% Substantive departures, each footnoted where it occurs:
%   - no. 6: Lambda is characterised by L Lambda = 1 - sum pi_{i,0} and
%     Lambda L = 1 - sum pi_{j,0}, which makes it a two-sided pseudo-inverse of
%     L and NOT the Hodge-theoretic Lambda (for which Lambda L acts by n-i on
%     P^i). The normalisation is what keeps everything over Z; the pages never
%     remark on it.
%   - p. 12: the margin writes Lambda in E^{(2)} where Lambda has degree -2;
%     the sign is missing on the leaf.
%   - p. 10: a second equation carries the label (6.7), duplicating p. 9's.
%   - p. 3: a reference « (1.2.) » points at what p. 2 labels (1.3) b).
%   - pp. 4-5: u_i and u_{2n-i} are written where the argument needs pi_i and
%     pi_{2n-i}.
%   - p. 43: « C(X) is of finite dimension over Q » uses C for the group of
%     cycles modulo numerical equivalence, not for the conjecture C(X) two
%     lines above; and « C^{n-1}_chi » should read C^{n-i}.
%   - pp. 43 and 47: the two letters are dated 4.1.67 and 6.1.1966 in his hand.
%     Both dates are transcribed; nothing on the leaves reconciles them.
%   - p. 47: the Cayley-Hamilton relation is written in u with the coefficients
%     of w; restated in w here.
%   - p. 50: the leaf reads « Alors A contient L », with L drawn as the
%     Lefschetz operator. Since zeta lies in A by hypothesis, the informative
%     conclusion is that A contains the pseudo-inverse Lambda; stated that way,
%     with the reading flagged.
%   - pp. 52-54: the step « zeta^{n-k} T_k est nul sur H^j pour j =/= k » does
%     not follow from T_k as the leaf defines it. Recorded as a gap, not
%     repaired.
%
% The conjectures' LABELS are the folder's own, of January 1967, and do not
% match the ones published a year later: here C(X) is the conjecture that the
% inverse isomorphisms of hard Lefschetz are algebraic, D(X) that homological
% and numerical equivalence agree, A(T) the Lefschetz-type statement about
% algebraic cycles on T, with A^0 its restriction to the critical dimension.
% Each is given its meaning where it first appears, deduced from the pages'
% own use; no dictionary with the published names is asserted.
%
% Announced and never established, recorded in the spine section: his no. 8
% (p. 17) opens on necessary conditions and breaks off; « Démonstration à
% fournir » (p. 12); the Question of p. 13 asking for a synthetic presentation
% of the relations between L, Lambda and the pi_i; the finiteness of the orbit
% on p. 28; the gap in the Lemma's induction (p. 52); and, on p. 46, the
% reformulation of the standard conjectures for singular varieties.
%
% Modern names are footnoted with whether the pages could have known them. The
% dating 1966-1967 is the archivists'.
\documentclass[11pt,a4paper]{article}
% Le préambule commun à toutes les transcriptions du fonds Grothendieck.
%
% Il tient en une idée : **l'incertitude se compose, elle ne se résout pas.**
% Une transcription qui lisse un mot douteux a détruit la seule chose pour
% laquelle on la faisait — un lecteur doit pouvoir savoir, à chaque endroit,
% ce qui était sur le feuillet et ce qui a été ajouté. Les six macros qui
% suivent sont donc l'appareil critique, et non de la décoration.
%
% Les mêmes commandes sont reconnues par `scripts/render.mjs`, qui produit la
% vue de lecture du site. Une macro ajoutée ici doit l'être aussi là-bas,
% faute de quoi le rendu échouera bruyamment — ce qui est voulu : un rendu
% silencieusement faux est pire qu'un rendu absent.

\ProvidesPackage{grothendieck}[2026/08/08 Transcriptions du fonds Grothendieck]

\usepackage{amsmath,amssymb,amsthm}
\usepackage{mathrsfs}
\usepackage{stmaryrd}
% Les diagrammes commutatifs sont partout dans ce fonds : c'est la langue
% même de la géométrie algébrique de Grothendieck, et une transcription qui
% les rendrait en prose ne transcrirait rien.
\usepackage{tikz-cd}
% Français par défaut — c'est la langue des feuillets — mais l'anglais est
% chargé aussi : la traduction et le résumé sont des documents anglais, et la
% typographie française (espaces avant les deux-points) y serait une faute.
% Ces documents basculent par \selectlanguage{english} en tête de corps.
\usepackage[english,french]{babel}
\usepackage{fontspec}
\usepackage{geometry}
\geometry{a4paper,margin=25mm,marginparwidth=28mm}
\usepackage{marginnote}
\usepackage{xcolor}
% ulem, not soul. `soul`'s \st and \ul are built on a character-by-character
% scan that XeLaTeX's font machinery breaks: the document fails with an
% undefined control sequence rather than degrading. ulem does the same two jobs
% and is Unicode-engine safe. `normalem` stops it hijacking \emph, which the
% transcriptions use for Grothendieck's own emphasis.
\usepackage[normalem]{ulem}
\usepackage{hyperref}
\hypersetup{colorlinks=true,linkcolor=black,urlcolor=black,pdfborder={0 0 0}}

% --- Métadonnées du lot -----------------------------------------------------
% Reprises telles quelles par le script de rendu, qui les affiche en tête de
% la vue de lecture. Elles fixent ce que le fichier prétend couvrir : sans
% elles, une transcription flotte sans point d'attache dans un fonds de seize
% mille feuillets.

\newcommand{\folder}[1]{\gdef\grothFolder{#1}}
\newcommand{\batch}[1]{\gdef\grothBatch{#1}}
\newcommand{\leaves}[2]{\gdef\grothFirst{#1}\gdef\grothLast{#2}}
\newcommand{\foldertitle}[1]{\gdef\grothFoldertitle{#1}}
\gdef\grothFolder{?}\gdef\grothBatch{?}\gdef\grothFirst{?}\gdef\grothLast{?}\gdef\grothFoldertitle{}

% --- Le résumé --------------------------------------------------------------
%
% Il ouvre la lecture modernisée, sous le titre « Résumé », et il est composé
% en petit. Ce n'est pas un ornement : c'est le seul passage du document qui
% ne soit pas des mathématiques, et un lecteur doit pouvoir le reconnaître
% avant de l'avoir lu — puis le sauter s'il connaît déjà le sujet, ou s'y
% arrêter s'il ne le connaît pas. Le filet qui le suit marque la reprise du
% corps.

\newenvironment{resume}
  {\small\setlength{\parskip}{0.45em}}
  {\par\medskip\hrule\medskip}

% --- Mots-clés ---------------------------------------------------------------
%
% En anglais, en clôture du résumé de la lecture modernisée : le vocabulaire
% moderne sous lequel chercher ce que les feuillets construisent. Le manifeste
% du site (`npm run manifest`) les extrait pour étiqueter la cote entière —
% cette ligne est la source unique des tags, il n'y a pas de fichier à côté.
\newcommand{\keywords}[1]{\par\smallskip\noindent
  {\footnotesize\textsc{Keywords} --- \textit{#1}}\par}

% --- L'appareil critique ----------------------------------------------------

% Le numéro de feuillet, en marge. C'est l'ancre qui permet de régler un
% désaccord sur une formule en regardant *un* feuillet plutôt qu'une chemise
% de six cents. Le site s'en sert aussi pour tourner les pages du fac-similé
% quand on fait défiler la transcription.
\newcommand{\leaf}[1]{%
  \marginnote{\footnotesize\color{gray!70}\textsf{#1}}%
  \label{leaf:#1}\ignorespaces}

% Illisible : on ne devine pas. Un mot inventé qui se lit comme les autres est
% le pire résultat possible.
\newcommand{\ill}{\textcolor{red!60!black}{[\,\ldots\,]}}

% Lecture douteuse : le mot est donné, et signalé comme incertain.
\newcommand{\uncertain}[1]{\uline{#1}}

% Ajout de l'éditeur — une lettre manquante, un mot sous-entendu.
\newcommand{\add}[1]{\textcolor{blue!50!black}{[#1]}}

% Biffé par Grothendieck lui-même. Ce qu'il a rayé fait partie du document :
% une rature dit souvent où la pensée a changé de direction.
\newcommand{\struck}[1]{\sout{#1}}

% Note du transcripteur, sur l'état matériel du feuillet.
\newcommand{\note}[1]{\par\smallskip{\small\color{gray!60!black}\textit{[#1]}}\par\smallskip}

% Note marginale de Grothendieck, distincte du corps du texte.
\newcommand{\marginal}[1]{\par\smallskip\noindent\hspace{1em}%
  {\small\color{gray!60!black}#1}\par\smallskip}

% --- Titre ------------------------------------------------------------------

\AtBeginDocument{%
  \begin{center}
    {\large\bfseries Fonds Alexandre Grothendieck --- cote n\textsuperscript{o}~\grothFolder}\\[2pt]
    {\itshape\grothFoldertitle}\\[4pt]
    {\small feuillets \grothFirst--\grothLast\ (lot \grothBatch)}\\[2pt]
    {\footnotesize\color{gray!60!black}Transcription établie d'après le fac-similé numérisé
     par l'Université de Montpellier.\\
     Les lectures incertaines sont soulignées, les illisibles notées [\,\ldots\,],
     les ajouts entre crochets.}
  \end{center}
  \vspace{1em}}

\endinput


\folder{16}
\pages{1}{55}
\dating{1966-1967}
\watermark{Édition de démonstration\\interprétation personnelle de l'œuvre}
\foldertitle{Formalisme algébrique des correspondances et algèbres de Poincaré-Hodge (cf conjectures standard) : notes manuscrites (s.d.), tapuscrit annoté (s.d.), lettres (1966-1967) — lecture modernisée du dossier entier}

\begin{document}

\section*{Résumé}

\begin{resume}
Ces feuillets posent une seule question, et la posent trois fois : on sait
monter, sait-on redescendre ?

Voici l'image. La cohomologie d'une variété projective lisse de dimension $n$
se range par étages, numérotés de $0$ à $2n$. La géométrie fournit un ascenseur
qui monte de deux étages : couper par une section hyperplane, c'est-à-dire
multiplier par la classe $\xi$ de cet hyperplan. On note $L$ cette opération.
Un théorème célèbre --- le \emph{théorème de Lefschetz difficile} --- dit que
cet ascenseur, pris $n-i$ fois de suite, réalise une bijection parfaite entre
l'étage $i$ et l'étage $2n-i$, symétrique du premier par rapport au milieu.
L'édifice est donc apparié étage par étage, et l'appariement est donné par la
géométrie.

Sait-on redescendre ? Il existe bien une application inverse, puisque la montée
est bijective ; la question est de savoir si elle est, elle aussi, donnée par
la géométrie --- autrement dit si elle est induite par une
\emph{correspondance algébrique}, une sous-variété du produit
$X \times X$. Rien ne l'assure : l'inverse d'une bijection construite
géométriquement n'a aucune raison d'être géométrique à son tour. C'est là
toute la difficulté, et c'est l'une des \emph{conjectures standard} sur les
cycles algébriques, aujourd'hui encore ouverte.

Le dossier attaque cette question par le bas. La première moitié (pages 2 à
17) oublie complètement la géométrie et ne garde qu'un anneau gradué
$\mathcal{E}$ d'opérateurs et un élément $L$ de degré $2$. Elle demande : que
faut-il exactement supposer pour que l'inverse $\Lambda$, et les projections
sur les étages, appartiennent déjà à $\mathcal{E}$ ? La réponse est un
théorème, et une surprise : il ne suffit pas de se donner $L$. Un exemple en
dimension $1$ le montre --- l'anneau engendré par $L$ y est si petit qu'il ne
sait rien des étages. Il faut donner les étages avec l'ascenseur. Une
seconde surprise vient trente pages plus loin : une algèbre de dimension six
où $L^{n}$ est non nul, et où pourtant l'appariement échoue au milieu. La
condition qui saute aux yeux n'est donc pas la bonne.

Entre les deux, quelques pages isolent la structure minimale qui porte tout
cela : une algèbre graduée finie, munie d'une dualité parfaite entre les
étages $i$ et $2n-i$ --- ce qu'il appelle une \emph{algèbre de Hodge}. C'est
l'ossature d'un anneau de cohomologie, dégagée de toute variété.

La seconde moitié redescend vers la géométrie, et le fait par la voie la plus
économique : une section hyperplane $Y \subset X$. Tout ce qui se passe sur
$X$ se laisse comparer à ce qui se passe sur $Y$, et les pages 32 à 39
établissent le formulaire complet de ces comparaisons. Il en sort une
réduction : la conjecture pour $X$ équivaut à la conjecture pour $Y$, plus un
énoncé portant sur une seule dimension, la dimension médiane. En descendant le
long d'une suite de sections hyperplanes emboîtées, tout se ramène ainsi à ce
qui arrive au milieu. Deux lettres à un correspondant qui rédigeait les notes
d'un exposé rassemblent ensuite l'état de la question, avec un argument
d'algèbre linéaire d'une simplicité désarmante : si un isomorphisme est
géométrique et si les polynômes caractéristiques des endomorphismes
géométriques sont rationnels, alors l'inverse est géométrique --- par
Cayley--Hamilton.

Deux phrases de ces lettres valent d'être lues pour elles-mêmes. L'une
s'étonne qu'un résultat de Hodge de 1938 soit « tombé dans l'oubli pendant
près de trente ans », et que personne n'y ait pris garde. L'autre, après avoir
énoncé le « yoga complet » --- il existerait une catégorie de motifs, et elle
serait semi-simple --- ajoute qu'il serait bon, « sans parler de motifs, et en
restant les pieds sur terre », de dire la même chose en termes d'anneaux de
correspondances. Le dossier entier tient dans cet écart : la vision qui
organise, et la vérification élémentaire qui, seule, se démontre.

\keywords{hard Lefschetz theorem, primitive decomposition, Lefschetz operator, sl2-representation, graded Artinian Gorenstein algebra, Poincaré duality algebra, strong Lefschetz property, standard conjectures on algebraic cycles, Lefschetz standard conjecture, Künneth projectors, algebraic correspondences, numerical equivalence, homological equivalence, Hodge index conjecture, hyperplane section, weak Lefschetz theorem, pure motives, semisimplicity, Cayley-Hamilton, Néron-Severi group, Picard scheme, Tate module}
\end{resume}

\pagerange{2}{2}
\section*{Le fil du dossier, et les conventions}

\emph{Les stations.}

\begin{enumerate}
  \item \emph{Pages 2 à 17}, « Sur certains anneaux gradués » : de l'algèbre
  pure. Un anneau gradué $\mathcal{E}$, un élément $L$ de degré $2$ ; à quelles
  conditions la graduation, l'opérateur $\Lambda$ et les projecteurs de la
  décomposition primitive sont-ils déterminés par ces données et internes à
  $\mathcal{E}$ ? Théorème du n° 6, anneau universel $\Phi_{0}$ du n° 7, et la
  réponse négative à la question « $L$ suffit-il ? ».
  \item \emph{Pages 18 à 21} : deux feuillets d'un autre texte, une page
  d'essais pour écrire $\Lambda$ comme polynôme en $L$, et une algèbre de
  dimension $6$ qui sépare la condition $L^{n} \neq 0$ de la propriété de
  Lefschetz.
  \item \emph{Pages 24 à 28} : les « algèbres de Hodge d'ordre 2 », ossature
  axiomatique de l'anneau de cohomologie d'une surface.
  \item \emph{Pages 32 à 39} : le formulaire de $L$ et $\Lambda$ le long d'une
  section hyperplane lisse $Y \subset X$, et le corollaire
  $(A(X \times X) \text{ et } C(Y)) \Rightarrow C(X)$.
  \item \emph{Pages 41 à 49} : la réduction $C(X) \iff [\,C(Y)$ et
  $A^{0}(X \times Y)\,]$, puis deux lettres à Coates.
  \item \emph{Pages 50 à 54} : le lemme qui fait entrer $\Lambda$ et les
  projecteurs dans un anneau de correspondances algébriques.
\end{enumerate}

\emph{Ce qui relie les stations, et qu'aucun feuillet ne dit.} Trois liens
courent d'un bout à l'autre du dossier sans être écrits nulle part.

Le premier est une identification de notation. Les $p^{X}_{i}$ du formulaire
(pages 32 à 39) sont les $\pi_{i,0}$ du premier cahier : les projections sur
les parties primitives. Les identités (4) et (5) du tapuscrit sont donc
exactement les relations (6.7) et (6.15) du n° 6, lues cette fois dans la
cohomologie d'une variété au lieu d'un anneau gradué abstrait.

Le deuxième est que le lemme des pages 50 à 54 est la moitié « existence » des
n° 5 et 6, transportée dans l'anneau des correspondances algébriques.
L'hypothèse de stabilité par composantes homogènes y joue le rôle que jouait,
page 7, le sur-anneau $\mathcal{F}$ dans lequel les projecteurs vivaient
d'avance.

Le troisième est que l'algèbre de dimension $6$ de la page 21 répond à la
liste de conditions nécessaires que le n° 8 ouvre page 17 et n'achève pas :
$L^{n} \neq 0$ est nécessaire, et n'est pas suffisant.

\emph{Les conventions.} Trois familles de projecteurs portent des lettres
voisines et sont distinguées ici une fois pour toutes :

\begin{itemize}
  \item $\pi_{i}$ (pages 2 à 17) : la projection du module gradué $M$ sur sa
  composante $M^{i}$ ;
  \item $\pi_{i,\alpha}$ (pages 8 à 17) : la projection sur le facteur
  élémentaire $L^{\alpha} P^{i-2\alpha}$ de la décomposition primitive ;
  \item $p^{X}_{i}$ (pages 32 à 39) : la projection sur la partie primitive
  $P^{i}(X)$ si $i \leqslant n$, et sur le sommet $L^{n-i'}P^{i'}$ avec
  $i = 2n-i'$ si $i > n$ ;
  \item $p_{i}$ (pages 50 à 54) : la projection de $H^{*}(X)$ sur $H^{i}(X)$.
\end{itemize}

$L$ désigne partout l'opérateur de Lefschetz, de degré $2$, et $\Lambda$ son
pseudo-inverse de degré $-2$ ; $P^{i}$ les parties primitives ;
$\varphi \colon Y \to X$ l'inclusion d'une section hyperplane lisse, avec
$\varphi^{*}$ et $\varphi_{*}$ ; $\xi_{X}$, $\xi_{Y}$ les classes de
polarisation ; $\chi$ l'indice de la théorie cohomologique choisie. Les
$\varphi$, $\psi$, $Q$, $\chi$ des pages 24 à 28 sont des formes bilinéaires
et n'ont aucun rapport avec l'inclusion $\varphi$ ni avec l'indice $\chi$ : la
collision est sur les feuillets et on la signale à chaque emploi.

\emph{Les lettres des conjectures} sont celles de janvier 1967 et non celles
qui seront publiées ensuite. Dans ce dossier $C(X)$ est la conjecture que les
isomorphismes inverses du théorème de Lefschetz difficile sont induits par des
correspondances algébriques ; $D(X)$ que l'équivalence homologique coïncide
avec l'équivalence numérique ; $A(T)$ l'énoncé de type Lefschetz portant sur
les cycles algébriques de $T$, et $A^{0}(T)$ sa restriction à la dimension
critique. Chacune reçoit son sens là où elle paraît.\footnote{Les
correspondances entre ces lettres et celles sous lesquelles les conjectures
standard ont été publiées ne sont pas établies ici : on lit les lettres du
dossier au sens que les feuillets leur donnent, et rien de plus. Une même
lettre y change d'ailleurs de sens d'une ligne à l'autre — voir la note sur la
page 43.}

\emph{Ce que le dossier annonce et n'établit pas.} Le n° 8 (page 17) ouvre la
liste des conditions nécessaires et s'arrête à la troisième. Une démonstration
est explicitement remise à plus tard page 12. La « Question » de la page 13,
qui demande une présentation synthétique des relations entre $L$, $\Lambda$ et
les $\pi_{i}$, reste sans réponse. L'esquisse de la page 28 sur la finitude
d'une orbite s'interrompt. L'induction du lemme (page 52) comporte une étape
que le feuillet ne justifie pas. Enfin la page 46 propose de reformuler les
conjectures standard pour les variétés singulières, et le dossier s'arrête là.

\pagerange{2}{7}
\section*{Quand la graduation est-elle déterminée par l'anneau ? (pages 2 à 7)}

Soit $k$ un anneau commutatif, $M$ un $k$-module gradué et $\mathcal{E}$ un
sous-anneau gradué de $\operatorname{End}(M)$. On note $\pi_{i}$ la projection
de $M$ sur $M^{i}$ et l'on suppose $\pi_{i} \in \mathcal{E}$ pour tout $i$.
La graduation de $\mathcal{E}$ s'exprime alors intrinsèquement en fonction des
seuls $\pi_{i}$ :
\[
\mathcal{E}^{(\nu)} = \sum_{i} \pi_{i+\nu}\, \mathcal{E}\, \pi_{i},
\qquad
u \in \mathcal{E}^{(\nu)} \iff u\,\pi_{i} = \pi_{i+\nu}\,u \ \ \text{pour tout } i,
\]
et la composante homogène de degré $\nu$ d'un endomorphisme $u$ est
$\sum_{i} \pi_{i+\nu}\,u\,\pi_{i}$.

Le mouvement propre de ces pages est de renverser cette description. On part
d'un anneau gradué $\mathcal{E}$ abstrait et d'une famille $(\pi_{i})_{i \in
\mathbb{Z}}$ d'éléments presque tous nuls, et l'on demande quand les formules
ci-dessus définissent bien la graduation. Il suffit qu'ils forment une famille
d'idempotents orthogonaux de somme $1$ :
\[
\pi_{i}^{2} = \pi_{i}, \qquad \pi_{i}\pi_{j} = 0 \ (i \neq j),
\qquad \sum_{i} \pi_{i} = 1 .
\]
On obtient alors du même coup la bigraduation
$\mathcal{E} = \bigoplus_{i,j} \pi_{j}\,\mathcal{E}\,\pi_{i}$, dont la
composante $(i,j)$ est $u \mapsto \pi_{j}\,u\,\pi_{i}$.\footnote{La page 3
renvoie à « (1.2.) » pour la formule que la page 2 étiquette (1.3) b) ; la
discordance est sur le feuillet et n'affecte pas l'énoncé.}

\subsection*{Unicité}

On fixe désormais un entier $n$ et l'on suppose $\pi_{i} = 0$ pour
$i \notin [0, 2n]$. La question devient : une telle famille est-elle unique ?

La première observation est que $\mathcal{E}^{(0)}$ est le commutant de la
sous-algèbre engendrée par les $\pi_{i}$, de sorte que les $\pi_{i}$
appartiennent au centre $\mathfrak{z}\mathcal{E}^{(0)}$, qui est un anneau
commutatif : ce sont des idempotents orthogonaux de somme $1$ de cet anneau.
Or, dans un anneau commutatif, un idempotent est connu dès qu'on connaît son
annulateur. Tout revient donc à déterminer les annulateurs, et c'est une
récurrence descendante.

Supposons connus les $\pi_{\alpha}$ pour $\alpha \notin [i, 2n-i]$, et posons
\[
\varphi_{i} = \sum_{\alpha \geqslant i} \pi_{\alpha},
\qquad
\psi_{2n-i} = \sum_{\alpha \leqslant 2n-i} \pi_{\alpha},
\]
qui sont donc connus. Un calcul immédiat --- tous les termes d'une somme
s'annulent sauf celui d'indice $i$ --- donne
\[
\psi_{2n-i}\ \mathcal{E}^{(2n-2i)}\ \varphi_{i}
= \pi_{2n-i}\,\mathcal{E}\,\pi_{i}
=: \mathcal{E}^{2n-i}_{i},
\]
de sorte que la composante bigraduée $\mathcal{E}^{2n-i}_{i}$ est elle aussi
connue. Les deux conditions qui referment la récurrence sont alors :
\[
\bigl( u \in \mathfrak{z}\mathcal{E}^{(0)},\ \ \mathcal{E}^{2n-i}_{i}\,u = 0 \bigr)
\implies \pi_{i}\,u = 0,
\qquad
\bigl( u \in \mathfrak{z}\mathcal{E}^{(0)},\ \ u\,\mathcal{E}^{2n-i}_{i} = 0 \bigr)
\implies \pi_{2n-i}\,u = 0 .
\]
Elles disent exactement que l'annulateur de $\pi_{i}$ dans
$\mathfrak{z}\mathcal{E}^{(0)}$ est l'idéal des $u$ tués à gauche par
$\mathcal{E}^{2n-i}_{i}$, et de même à droite pour $\pi_{2n-i}$ ; d'où
l'unicité.\footnote{Les feuillets 4 et 5 écrivent $u_{i}$ et $u_{2n-i}$ là où
l'argument porte sur $\pi_{i}$ et $\pi_{2n-i}$ ; les glyphes ne portent pas la
barre du $\pi$ et la transcription les laisse tels quels. On a rétabli ici ce
que l'énoncé demande.}

\subsection*{Une condition suffisante, puis interne}

Quand $\mathcal{E}$ agit sur un module $M$, il suffit pour cela que
l'intersection des noyaux des $u \colon M^{i} \to M^{2n-i}$, $u$ parcourant
$\mathcal{E}^{2n-i}_{i}$, soit nulle, et que la somme de leurs images soit
$M^{2n-i}$ ; ou, plus simplement, qu'il existe dans $\mathcal{E}^{(2n-2i)}$ un
élément induisant un isomorphisme $M^{i} \xrightarrow{\ \sim\ } M^{2n-i}$.

Cette dernière forme se transforme en une condition qui ne parle plus de $M$
du tout, et c'est le pas décisif de ces pages :
\[
(4.4)_{i} \qquad
\exists\, v_{i} \in \mathcal{E}^{(2n-2i)},\ w_{i} \in \mathcal{E}^{-(2n-2i)}
\quad\text{tels que}\quad
(w_{i}v_{i} - 1)\,\pi_{i} = 0, \quad (v_{i}w_{i} - 1)\,\pi_{2n-i} = 0 .
\]
On vérifie sans aucune hypothèse de réalisation que $(4.4)_{i}$ entraîne les
deux conditions d'unicité : si $\mathcal{E}^{2n-i}_{i}\,u = 0$ alors
$v_{i}u = 0$, donc $w_{i}v_{i}\pi_{i}u = 0$, c'est-à-dire $\pi_{i}u = 0$
puisque $w_{i}v_{i}\pi_{i} = \pi_{i}$.

\subsection*{Existence}

Le n° 5, écrit après coup et inséré sur un feuillet « 5 bis », tire de la même
condition l'existence. Supposons $\mathcal{E}$ contenu dans un sur-anneau
gradué $\mathcal{F}$ dans lequel les $\pi_{i}$ vivent déjà, et supposons
$(4.4)_{i}$ satisfaite avec $v_{i}, w_{i} \in \mathcal{E}$. Alors les
$\pi_{i}$ sont eux-mêmes dans $\mathcal{E}$.

L'argument est d'une brièveté exemplaire. Par récurrence, $\varphi_{i}$ et
$\psi_{2n-i}$ sont dans $\mathcal{E}$ ; on pose
\[
v'_{i} = \psi_{2n-i}\,v_{i}\,\varphi_{i} = \pi_{2n-i}\,v_{i}\,\pi_{i} \in \mathcal{E},
\qquad
w'_{i} = \varphi_{i}\,w_{i}\,\psi_{2n-i} = \pi_{i}\,w_{i}\,\pi_{2n-i} \in \mathcal{E},
\]
et l'on calcule, en faisant passer $w_{i}$ à travers $\pi_{2n-i}$ et en
utilisant $w_{i}v_{i}\pi_{i} = \pi_{i}$ :
\[
w'_{i}v'_{i} = \pi_{i}\,w_{i}\,\pi_{2n-i}\,v_{i}\,\pi_{i} = \pi_{i}^{3} = \pi_{i},
\qquad v'_{i}w'_{i} = \pi_{2n-i} .
\]
Les deux projecteurs sont donc des produits d'éléments de $\mathcal{E}$.

\pagerange{8}{13}
\section*{La condition de Lefschetz, la décomposition primitive, et $\Lambda$ (pages 8 à 13)}

Le n° 6 spécialise tout ce qui précède au seul cas qui intéresse la géométrie :
celui où l'élément $v_{i}$ est une puissance d'un même opérateur.

Soit donc $L \in \mathcal{E}^{(2)}$, et considérons la
\emph{condition de Lefschetz}
\[
(6.1)_{i} \qquad L^{n-i} \colon M^{i} \longrightarrow M^{2n-i}
\quad\text{est un isomorphisme,}
\]
complétée par l'exigence que l'inverse appartienne à
$\mathcal{E}^{-(2n-2i)}$. Sous la forme interne, cela s'écrit
\[
(6.3)_{i} \qquad
\exists\, w_{i} \in \mathcal{E}^{-(2n-2i)} \ :\quad
(w_{i}L^{n-i} - 1)\,\pi_{i} = 0, \quad (L^{n-i}w_{i} - 1)\,\pi_{2n-i} = 0,
\]
et c'est un cas particulier de $(4.4)_{i}$ : l'unicité et l'appartenance des
$\pi_{i}$ à $\mathcal{E}$ sont donc acquises.

La condition entraîne la \emph{décomposition primitive} de $M$ :
\[
M^{i} \simeq P^{i} \oplus L P^{i-2} \oplus L^{2}P^{i-4} \oplus \cdots
\quad (i \leqslant n),
\qquad
M^{2n-i} \simeq L^{n-i}P^{i} \oplus L^{n-i+1}P^{i-2} \oplus \cdots,
\]
et l'on note $\pi_{i,\alpha}$ la projection de $M^{i}$ sur
$L^{\alpha}P^{i-2\alpha}$, $\pi_{2n-i,\alpha}$ celle de $M^{2n-i}$ sur
$L^{n-i+\alpha}P^{i-2\alpha}$.\footnote{C'est la décomposition de Lefschetz.
Le dossier ne la démontre pas et ne la nomme pas : elle est prise comme
conséquence connue de $(6.1)$.}

\subsection*{Le théorème du n° 6}

Le feuillet 9 énonce un théorème sous forme d'une chaîne de conditions
équivalentes. La liste est raturée et renumérotée trois fois, et seule sa
charpente se lit ; ce qui est sûr, et ce que la démonstration des pages 10 à
12 établit effectivement, est l'équivalence des trois énoncés
suivants.\footnote{Les énoncés intermédiaires de la liste manuscrite sont
biffés et réécrits au point que leur libellé ne se reconstitue pas ; on ne
restitue ici que ce que la démonstration utilise et prouve.}

\begin{enumerate}
  \item Les conditions de Lefschetz $(6.1)_{i}$ et $(6.2)_{i}$ valent pour tout
  $i$, et les $\pi_{i}$ sont dans $\mathcal{E}$.
  \item Il existe $\Lambda \in \mathcal{E}^{(-2)}$ satisfaisant les deux
  relations $(6.7)$ ci-dessous.
  \item Tous les projecteurs $\pi_{i,\alpha}$ de la décomposition primitive
  sont dans $\mathcal{E}$.
\end{enumerate}

Les deux relations qui déterminent $\Lambda$ sont
\[
(6.7) \qquad
L\Lambda = 1 - \sum_{0 \leqslant i \leqslant n} \pi_{i,0},
\qquad
\Lambda L = 1 - \sum_{n \leqslant i \leqslant 2n} \pi_{i,0} .
\]

Il faut s'arrêter sur elles, car elles ne disent pas ce qu'on attend. Elles
font de $\Lambda$ un \emph{pseudo-inverse à deux côtés} de $L$ : $\Lambda L$
vaut l'identité partout sauf sur les parties primitives situées au bas des
chaînes, et $L\Lambda$ vaut l'identité partout sauf sur les sommets. Ce n'est
pas l'opérateur $\Lambda$ de la théorie de Hodge, l'adjoint de $L$ pour la
métrique, lequel satisfait $[L,\Lambda] = H$ et agit sur $L^{k}P^{i}$ par des
scalaires entiers non triviaux --- sur une classe primitive de degré $i$ on a
$\Lambda L x = (n-i)\,x$ et non $x$.\footnote{Le choix n'est pas anodin et les
feuillets ne le commentent pas : c'est précisément la normalisation en
pseudo-inverse qui permet à tout le n° 6 de rester sur $\mathbb{Z}$, sans
introduire les dénominateurs qu'une représentation de $\mathfrak{sl}_{2}$
impose. En contrepartie $L$ et $\Lambda$ n'engendrent pas ici une algèbre de
Lie, et la relation de commutation qui identifierait la graduation ne vaut
plus --- ce dont la « Question » de la page 13 est le symptôme.}

\subsection*{La démonstration}

Elle procède par une double récurrence, descendante sur $i$ et, à $i$ fixé,
descendante sur $\alpha$. Supposons connus les $\pi_{j,\alpha}$ pour
$j \notin [i, 2n-i]$, et les $\pi_{i,\alpha'}$, $\pi_{2n-i,\alpha'}$ pour
$\alpha' > \alpha$. On pose
\[
q_{i,\alpha} = \sum_{\alpha' \leqslant \alpha} \pi_{i,\alpha'}
= 1 - \sum_{\alpha' > \alpha} \pi_{i,\alpha'} \ \in \mathcal{E} .
\]
L'opérateur $L^{n-i+\alpha}q_{i,\alpha}$ tue alors tous les facteurs
$L^{\alpha'}P^{i-2\alpha'}$ avec $\alpha' \neq \alpha$, et induit sur le
facteur restant l'isomorphisme
$L^{n-i+2\alpha} \colon P^{i-2\alpha} \xrightarrow{\ \sim\ }
L^{n-i+2\alpha}P^{i-2\alpha}$. En le composant avec l'inverse $w_{i-2\alpha}$
et en remontant par $L^{\alpha}$, on obtient la formule qui referme la
récurrence :
\[
\pi_{i,\alpha}
= L^{\alpha}\,w_{i-2\alpha}\,L^{n-i+\alpha}
  \Bigl( 1 - \sum_{\alpha' > \alpha} \pi_{i,\alpha'} \Bigr) \ \in \mathcal{E},
\qquad
\pi_{2n-i,\alpha} = L^{n-i}\,\pi_{i,\alpha}\,w_{i} .
\]

Il reste $\Lambda$. Comme $\Lambda = \sum_{j,\alpha} \Lambda\,\pi_{j,\alpha}$,
il suffit de traiter chaque facteur, et le calcul se fait en trois lignes :
\[
\Lambda\,\pi_{i,0} = 0,
\qquad
\Lambda\,\pi_{i,\alpha} = w_{i-2\alpha}\,L^{n-i+\alpha}\,\pi_{i,\alpha}
\ \ (\alpha \geqslant 1),
\qquad
\Lambda\,\pi_{2n-i,\alpha} = L^{n-i+\alpha-1}\,w_{i-2\alpha}\,L^{\alpha}\,\pi_{2n-i,\alpha} .
\]
La première ligne est la traduction de $(6.7)$ : $\Lambda$ annule les
primitifs du bas.\footnote{Le feuillet 12 porte en marge « $\Lambda \in
\mathcal{E}^{(2)}$ » ; $\Lambda$ est de degré $-2$ partout ailleurs et le
signe manque simplement sur la page. Le même feuillet marque « Démonstration à
fournir » pour l'implication qui refermerait la chaîne, et le dossier ne
revient pas dessus. Enfin la page 10 réutilise l'étiquette (6.7), déjà portée
page 9 ; il ne corrige pas.}

Une \emph{Remarque} reprend alors toute la caractérisation en $L$ et $\Lambda$
seuls, sans plus nommer les $w_{i}$ :
\[
(\Lambda^{n-i}L^{n-i} - 1)\,\pi_{i} = 0,
\qquad
(L^{n-i}\Lambda^{n-i} - 1)\,\pi_{2n-i} = 0,
\]
\[
\pi_{i,\alpha} = L^{\alpha}\,\Lambda^{n-i+2\alpha}\,L^{n-i+\alpha}
\Bigl(1 - \sum_{\alpha' > \alpha}\pi_{i,\alpha'}\Bigr),
\qquad
\pi_{2n-i,\alpha} = L^{n-i}\,\pi_{i,\alpha}\,\Lambda^{n-i} .
\]

\pagerange{13}{17}
\section*{L'anneau universel $\Phi_{0}$, et ce que $L$ seul ne détermine pas (pages 13 à 17)}

Le n° 7 construit le modèle universel. On part d'un $\mathbb{Z}$-module libre
de rang fini $M_{0}$, d'une forme $\xi_{0} \in \Lambda^{2}M_{0}$ et de sa duale
$\xi_{0}^{*}$, on prend pour module gradué l'algèbre extérieure
$M_{0}^{*} = \Lambda M_{0}$, et l'on y fait agir
\[
L_{0}\,x = \xi_{0} \wedge x,
\qquad
\Lambda_{0}\,x = \xi_{0}^{*} \lrcorner\, x .
\]
On pose enfin
\[
\Phi_{0} = \mathbb{Z}[L_{0}, \Lambda_{0}] \subset \operatorname{End}(\Lambda M_{0}) .
\]
Cet anneau gradué est engendré par les éléments
\[
\varphi^{\beta,\alpha}_{i} =
\begin{cases}
L_{0}^{\beta-\alpha}\,\pi_{i,\alpha} = \pi_{j,\beta}\,L_{0}^{\beta-\alpha}
 & \text{si } \alpha \leqslant \beta, \\[2pt]
\Lambda_{0}^{\alpha-\beta}\,\pi_{i,\alpha} = \pi_{j,\beta}\,\Lambda_{0}^{\alpha-\beta}
 & \text{si } \alpha \geqslant \beta,
\end{cases}
\qquad 0 \leqslant i \leqslant n, \quad 0 \leqslant \alpha, \beta \leqslant 2n-i,
\]
dont la table de multiplication ne dépend pas du choix de
$(M_{0}, \xi_{0})$.\footnote{Le début du n° 7 est écrit très vite et sa syntaxe
se perd ; les données $M_{0}$, $\xi_{0}$, $\xi_{0}^{*}$ et l'algèbre extérieure
sont ce que la suite du calcul exige, et la transcription le signale. Un
premier jeu d'indices, plus large, est biffé sur le feuillet.} L'élément
$\varphi^{\beta,\alpha}_{i}$ envoie le facteur $L^{\alpha}P^{i-2\alpha}$ sur le
facteur $L^{\beta}P^{i-2\alpha}$ et tue les autres : $\Phi_{0}$ est donc
l'anneau de toutes les façons de circuler dans les chaînes de la décomposition
primitive.

\subsection*{La proposition, et son corollaire}

La \emph{Proposition} du feuillet 14 dit alors que, pour $\mathcal{E}$ et
$L \in \mathcal{E}^{(2)}$ donnés, les conditions suivantes sont équivalentes :

\begin{enumerate}
  \item $L$ satisfait les conditions du théorème du n° 6 ;
  \item il existe dans $\mathcal{E}$ un $\Lambda$ les satisfaisant ;
  \item il existe un homomorphisme d'anneaux gradués
  $\varphi \colon \Phi_{0} \to \mathcal{E}$ envoyant $L_{0}$ sur $L$ et les
  $\pi_{i}$ sur les $\pi_{i}$.
\end{enumerate}

Sous ces conditions l'homomorphisme est unique, et il transporte
automatiquement $\Lambda_{0}$ sur $\Lambda$, les $\pi_{i,\alpha}$ sur les
$\pi_{i,\alpha}$, les $\varphi^{\beta,\alpha}_{i}$ sur leurs homonymes. Le
\emph{Corollaire} du feuillet 16 en tire la version qui ne suppose plus de
module : si un tel homomorphisme existe, alors les $\pi_{i}$ sont déterminés
par la seule structure d'anneau gradué de $\mathcal{E}$.

\subsection*{La réponse est « Non »}

Suit une \emph{Question}, qu'il pose et à laquelle il répond en marge d'un seul
mot : un homomorphisme $\Phi_{0} \to \mathcal{E}$ avec $u(L_{0}) = L$ suffit-il,
sans qu'on impose les $\pi_{i}$ ?

\emph{Non} --- et le contre-exemple tient en une ligne. Prenons
$n = 1$, $2n = 2$. L'anneau $\Phi_{0}$ contient $\pi_{0}$, $\pi_{1}$,
$\pi_{2}$, $L$ et $\Lambda$ ; mais le sous-anneau engendré par $L$ seul est
$\mathbb{Z} + \mathbb{Z}L$, et rien d'autre. Il ne contient aucun des trois
projecteurs. « Il faut donc $L$ \emph{et} les $\pi_{i}$. »

C'est la conclusion du premier cahier, et elle est plus qu'une précaution
technique : elle dit que l'opérateur de Lefschetz ne porte pas en lui la
graduation qu'il relie. Dans la théorie de Hodge, c'est le commutateur
$[L, \Lambda]$ qui restitue le degré ; ici, avec un $\Lambda$ normalisé en
pseudo-inverse, ce commutateur vaut
$\sum_{i \geqslant n}\pi_{i,0} - \sum_{i \leqslant n}\pi_{i,0}$ et ne retient
que les extrémités des chaînes. La graduation doit donc être fournie
séparément.\footnote{Ce calcul est une conséquence immédiate des deux
relations $(6.7)$ ; le dossier ne l'écrit pas et ne mentionne jamais
$\mathfrak{sl}_{2}$.}

\subsection*{Le n° 8, resté ouvert}

Le dernier feuillet du cahier ouvre un n° 8 sur les conditions
\emph{nécessaires} pour qu'un $L$ satisfasse les conditions de Lefschetz ---
qu'il écrit ici en toutes lettres et souligne --- et s'arrête après trois :
\[
\mathcal{E}^{(2n)} \neq 0 \iff \pi_{2n} \neq 0
\implies L^{n} \neq 0 \implies \Lambda^{n} \neq 0 .
\]
La suite n'est pas sur ces feuillets. Mais la page 21 y répond par la négative,
et c'est le premier des liens que ce dossier ne fait pas lui-même.

\pagerange{18}{21}
\section*{$\Lambda$ comme polynôme en $L$, et une algèbre où la condition $L^{n} \neq 0$ ne suffit pas (pages 18 à 21)}

\subsection*{Un essai (page 19)}

Un feuillet de calculs, sans phrase liante, poursuit la Question de la page 13
sous la forme la plus concrète : peut-on écrire $\Lambda$ comme un polynôme en
$L$ ? Deux lignes sont commencées et laissées sans second membre,
$\Lambda = P(L$ et $\Lambda = \Omega(L)\Delta(L)$.

Il essaie alors sur un exemple à deux générateurs $L, L'$ de degré $2$, en
écrivant une base graduée $L, L' \mid L^{2}, LL' \mid L^{3}$ et les relations
qui la ferment :
\[
L'^{2} = LL', \qquad L^{3} = L^{2}L', \qquad
L'^{3} = L'\,L'^{2} = L^{2}L' = L^{3},
\]
puis compare l'anneau obtenu à $\mathbb{Z}[L]/(L^{4})$.\footnote{La dernière
ligne de la page n'est pas lisible jusqu'au bout et la comparaison n'aboutit
pas sur le feuillet. Les deux feuillets qui l'encadrent, pages 18 et 20, ne
relèvent pas de ce dossier : ce sont deux pages d'un autre texte de sa main,
numérotées 42 et 45 dans leur pagination propre, sur le schéma de Picard
$\underline{\mathrm{Pic}}_{U}$, la suite exacte
$0 \to \mathrm{Pic}(U)^{0} \to \mathrm{Pic}(U) \to \mathrm{NS}(U) \to 0$, la
définition négative de la matrice d'intersection des composantes de la fibre
d'une résolution --- d'où l'autodualité de $\mathrm{NS}(U)$ à valeurs dans
$\mathbb{Q}/\mathbb{Z}$ --- et la suite exacte $\ell$-adique faisant intervenir
le module de Tate. Elles portent des corrections de sa main, ce qui est la
seule raison pour laquelle la transcription les garde.}

\subsection*{Le contre-exemple (page 21)}

La page 21 est un feuillet au crayon, biffé aux trois quarts ; la prose se perd
mais les formules portent, et ce qu'elles portent est décisif.

Il commence par vérifier une normalisation : si $L'^{n} = \alpha L_{0}^{n}$ et
$\Lambda'^{n} = \beta \Lambda_{0}^{n}$, alors
\[
L_{0}^{n} = L_{0}^{n}\pi_{0} = L_{0}^{n}\Lambda'^{n}L'^{n}
= L_{0}^{n}(\beta\Lambda_{0}^{n})(\alpha L_{0}^{n})
= \alpha\beta\,L_{0}^{n}\pi_{0} = \alpha\beta\,L_{0}^{n},
\]
d'où $\alpha\beta = 1$ : deux opérateurs de Lefschetz proportionnels en degré
maximal le sont par une unité.

Puis vient l'algèbre. Prenons la $k$-algèbre graduée de dimension $6$,
engendrée en degré $2$ par deux éléments $L$ et $L'$ soumis aux relations
\[
L^{4} = 0, \qquad L'^{2} - LL' = 0, \qquad L^{2}L' - L^{3} = 0, \qquad L^{3}L' = 0 .
\]
Une base graduée en est
\[
\underbrace{1}_{\deg 0} \ \Big\| \ \underbrace{L,\ L'}_{\deg 2}
\ \Big\| \ \underbrace{L^{2},\ LL'}_{\deg 4} \ \Big\| \ \underbrace{L^{3}}_{\deg 6},
\]
donc $2n = 6$ et $n = 3$. Alors :

\begin{itemize}
  \item $L'^{3} = L'\cdot L'^{2} = L'\cdot LL' = L\cdot L'^{2} = L\cdot LL'
  = L^{2}L' = L^{3} \neq 0$ ;
  \item $L$ est de Lefschetz : $L^{3}$ est un isomorphisme du degré $0$ sur le
  degré $6$, et $L$ un isomorphisme du degré $2$ sur le degré $4$ ;
  \item $L'$ ne l'est pas : $L' \cdot L = LL'$ et $L' \cdot L' = L'^{2} = LL'$,
  de sorte que $L'$ envoie tout le degré $2$ sur la seule droite engendrée par
  $LL'$, et n'est donc pas injectif en degré $2$.
\end{itemize}

Ainsi $L'^{n} = L^{n} \neq 0$ sans que $L'$ soit de Lefschetz. La chaîne de
conditions nécessaires ouverte page 17 --- $\pi_{2n} \neq 0 \Rightarrow L^{n}
\neq 0$ --- n'est donc pas réversible, et l'échec se produit exactement à
l'étage médian.\footnote{C'est aujourd'hui la distinction entre la
non-annulation de la puissance maximale et la \emph{propriété de Lefschetz
forte} d'une algèbre graduée artinienne de Gorenstein ; le vocabulaire est
postérieur de plus de dix ans à ces feuillets, qui ne le connaissent
évidemment pas et n'énoncent pas la distinction sous forme générale. Le calcul
ci-dessus est reconstitué à partir des seules relations que la page écrit ; la
phrase qui l'introduit y est en grande partie illisible.}

\pagerange{24}{28}
\section*{Les algèbres de Hodge d'ordre 2 (pages 24 à 28)}

Ces trois feuillets dégagent l'ossature d'un anneau de cohomologie de surface,
sans surface.\footnote{Le titre « Algèbres de Hodge » est de sa main : il
figure seul en tête d'un feuillet par ailleurs vierge qui sert de chemise. Les
lettres $\varphi$, $\psi$, $Q$, $\chi$ y désignent des formes bilinéaires et
n'ont rien à voir avec l'inclusion $\varphi$ ni avec l'indice $\chi$ des pages
suivantes.}

Une \emph{algèbre de Hodge d'ordre 2} est une algèbre graduée
anticommutative sur $k$, concentrée en degrés $0$ à $4$ :
\[
\begin{array}{ccccc}
H^{0} & H^{1} & H^{2} & H^{3} & H^{4} \\
k & V & k\xi \oplus W & \check{V} & k\xi^{2}
\end{array}
\]
dont la multiplication est entièrement déterminée par trois données : une
forme alternée $\varphi$ sur $V$, une application $\psi \colon \Lambda^{2}V \to
W$, et une forme bilinéaire symétrique $Q$ sur $W$. Elles se lisent sur les
deux applications
\[
\Lambda^{2}V \xrightarrow{\ \alpha\ } k\xi \oplus W,
\quad \alpha(x \wedge y) = \varphi(x,y)\,\xi + \psi(x,y),
\qquad
\mathrm{Sym}^{2}W \xrightarrow{\ \beta\ } k\xi^{2},
\quad \beta(xy) = Q(x,y)\,\xi^{2} .
\]
Les produits restants sont alors forcés : $x\xi = \tilde\varphi(x)$ où
$\tilde\varphi \colon V \to \check{V}$ est l'application associée à $\varphi$ ;
$xw = \rho(w,x) \colon y \mapsto Q(\psi(y,x), w)$ ; $x y' = \langle x, y'
\rangle\,\xi^{2}$ pour $x \in V$, $y' \in \check{V}$ ; $\xi w = 0$ ;
$ww' = Q(w,w')\,\xi^{2}$.

Quels que soient $\varphi$, $\psi$, $Q$, on obtient ainsi une loi
anticommutative unifère. Trois conditions la qualifient, et c'est le point de
ces pages qu'elles se lisent chacune sur une seule des trois données :

\begin{itemize}
  \item \emph{associativité} : $(xy)(tz) = (tx)(yz)$, soit
  \[
  \varphi(x,y)\varphi(t,z) + Q(\psi(x,y), \psi(t,z))
  = \varphi(t,x)\varphi(y,z) + Q(\psi(t,x), \psi(y,z)) ;
  \]
  \item \emph{condition de Lefschetz} : $\xi^{2} \colon H^{0} \to H^{4}$ est
  toujours un isomorphisme, et $\xi \colon H^{1} \to H^{3}$ en est un si et
  seulement si $\varphi$ est non dégénérée ;
  \item \emph{condition de Poincaré} : les accouplements $H^{0} \times H^{4}$
  et $H^{1} \times H^{3}$ vers $H^{4}$ sont toujours non dégénérés, et
  $H^{2} \times H^{2} \to H^{4}$ l'est si et seulement si $Q$ est non
  dégénérée.
\end{itemize}

\subsection*{Le cas dégénéré, et la reformulation}

Le feuillet 26 examine le cas où $(xy)(zt) = 0$, c'est-à-dire où l'image de
$\alpha$ est isotrope. Plus généralement, il introduit la forme
quadrilinéaire alternée $\chi$ définie par
$x\,y\,z\,t = \chi(x,y,z,t)\,\xi^{2}$, et observe que la forme $Q \oplus 1$
induite sur $\Lambda^{2}V$ par $\alpha$ se factorise à travers $\chi$. Toute
l'algèbre se ramène alors à la donnée de

\begin{itemize}
  \item un espace $V$ de dimension finie sur $k$ ;
  \item une forme $\chi \colon \Lambda^{4}V \to k$, d'où une forme symétrique
  $\chi'$ sur $\Lambda^{2}V$ ;
  \item un espace $\tilde{W}$ muni d'une forme bilinéaire symétrique
  $\tilde{Q}$ et d'une application $\alpha \colon \Lambda^{2}V \to \tilde{W}$
  compatible avec $\chi'$ ;
  \item un élément $\xi \in \tilde{W}$ avec $Q(\xi,\xi) = 1$,
\end{itemize}

les conditions de non-dégénérescence devenant : $Q$ non dégénérée, et
$\varphi(x,y) = Q(\alpha(x \wedge y), \xi)$ non dégénérée. Le gain est que la
polarisation $\xi$ devient un \emph{paramètre} de la structure, et non plus une
partie de sa définition --- ce qui est exactement la question que les lettres
de janvier poseront pour les conjectures.

L'exemple du feuillet 28 prend $\chi = \varphi^{2}$, $\alpha = \mathrm{id}$,
$\tilde{W} = \Lambda^{2}V$ ; $\check\varphi^{(2)}$ définit alors un
endomorphisme $u$ de $\Lambda^{2}V$, et il se demande si l'orbite de $\xi$ et
de $\check\varphi$ sous $u$ peut être infinie lorsque $\dim V \geqslant 4$.
L'esquisse s'interrompt.\footnote{La syntaxe de ce dernier passage ne se
reconstitue pas ; la question de finitude reste posée et sans réponse sur le
feuillet.}

\pagerange{32}{39}
\section*{Le formulaire de $L$ et $\Lambda$ le long d'une section hyperplane (pages 32 à 39)}

On passe ici à la géométrie. $X$ est une variété projective lisse connexe de
dimension $n$, polarisée par $\xi_{X} \in H^{2}(X)$ ; $L_{X}$ est le cup-produit
par $\xi_{X}$ et $\Lambda_{X}$ le pseudo-inverse de degré $-2$ ;
$p^{X}_{i}$ désigne la projection sur la partie primitive $P^{i}(X)$.
$Y \subset X$ est une section hyperplane lisse, $\varphi \colon Y \to X$
l'inclusion, $\xi_{Y}$ la polarisation induite.\footnote{Ces $p^{X}_{i}$ sont
les $\pi_{i,0}$ du premier cahier, étendus aux indices $i > n$ comme
projections sur les sommets $L^{n-i'}P^{i'}$, $i = 2n - i'$. Le dossier ne fait
nulle part le rapprochement, mais il est forcé : c'est ce qui rend les
identités (4) et (5) ci-dessous identiques aux relations $(6.7)$ de la page 9.
Le tapuscrit écrit l'inclusion avec un O barré, rendu ici $\varphi$ comme le
fait sa propre main dans la suite manuscrite.}

\subsection*{Les identités de base}

Des descriptions
\[
\operatorname{Im} L_{X} = \bigcap_{0 \leqslant i \leqslant n} \operatorname{Ker} p^{X}_{i},
\qquad
\operatorname{Ker} L_{X} = \sum_{n \leqslant i \leqslant 2n} \operatorname{Im} p^{X}_{i},
\]
et de leurs transposées pour $\Lambda_{X}$, on tire immédiatement
$p^{X}_{i}L_{X} = 0$ pour $i \leqslant n$, $L_{X}p^{X}_{i} = 0$ pour
$i \geqslant n$, les relations symétriques pour $\Lambda_{X}$, et en
particulier l'annulation complète en degré médian :
\[
p^{X}_{n}L_{X} = L_{X}p^{X}_{n} = p^{X}_{n}\Lambda_{X} = \Lambda_{X}p^{X}_{n} = 0 .
\]
Viennent ensuite les deux relations fondamentales, qui sont $(6.7)$ :
\[
\Lambda_{X}L_{X} = \mathrm{id}_{X} - \sum_{n \leqslant i \leqslant 2n} p^{X}_{i},
\qquad
L_{X}\Lambda_{X} = \mathrm{id}_{X} - \sum_{0 \leqslant i \leqslant n} p^{X}_{i} .
\]

\subsection*{Ce que la section hyperplane apporte}

Les formules de projection donnent
\[
\varphi^{*}L_{X} = L_{Y}\varphi^{*}, \qquad
L_{X}\varphi_{*} = \varphi_{*}L_{Y}, \qquad
\varphi_{*}\varphi^{*} = L_{X}, \qquad \varphi^{*}\varphi_{*} = L_{Y} ,
\]
la dernière se vérifiant en composant avec $\varphi^{*}$ et en constatant que
les deux membres deviennent $L_{X}^{2}$. Il en résulte la formule clef du
tapuscrit :
\[
\varphi_{*}\Lambda_{Y}\varphi^{*}
= \mathrm{id}_{X} - \sum_{0 \leqslant i \leqslant 2n} p^{X}_{i}
= (\Lambda_{X}L_{X} + L_{X}\Lambda_{X}) - \mathrm{id}_{X} + p^{X}_{n},
\]
la seconde égalité s'obtenant en additionnant les deux relations précédentes,
où $p^{X}_{n}$ est compté deux fois. Il pose alors
\[
T_{X} = \mathrm{id}_{X} + \varphi_{*}\Lambda_{Y}\varphi^{*},
\qquad\text{de sorte que}\qquad
\Lambda_{X}L_{X} + L_{X}\Lambda_{X} = T_{X} - p^{X}_{n},
\]
d'où l'identité de degré supérieur
\[
L_{X}^{2}\Lambda_{X} + 2\,L_{X}\Lambda_{X}L_{X} + \Lambda_{X}L_{X}^{2}
= L_{X}T_{X} + T_{X}L_{X} .
\]

Le \emph{corollaire} que le tapuscrit en tire est le premier énoncé du dossier
sur les conjectures :
\[
\bigl( A(X \times X) \ \text{et}\ C(Y) \bigr) \implies C(X),
\]
et une marge précise qu'on peut y remplacer $A(X \times X)$ par
$A^{0}(X \times X)$, la condition $A$ restreinte aux correspondances en
dimension critique.

\subsection*{La comparaison des décompositions primitives}

Les feuillets manuscrits comparent ensuite terme à terme les décompositions de
$X$ et de $Y$. Le théorème de Lefschetz faible donne
$P^{i}(X) \xrightarrow{\ \sim\ } P^{i}(Y)$ pour $i \leqslant n-2$, une
injection $P^{n-1}(X) \hookrightarrow P^{n-1}(Y)$, et
$P^{i}(X) = \operatorname{Ker}\varphi^{i}$ pour $i \geqslant n$, le tout se
lisant sur le carré

\begin{tikzcd}[column sep=large]
H^{i}(X) \arrow[r, "L_{X}^{n-i+1}"] \arrow[d, "\varphi^{*}"'] & H^{2n-i+2}(X) \\
H^{i}(Y) \arrow[r, "L_{Y}^{n-i}"'] & H^{2(n-1)-i}(Y) \arrow[u, "\varphi_{*}"']
\end{tikzcd}

pour $i \leqslant n-2$, condition sous laquelle $2(n-1) - i \geqslant n$. Tout
l'écart entre les deux variétés se concentre alors dans un seul module, le
conoyau en degré médian :
\[
E^{n-1} = H^{n-1}(Y)/\operatorname{Im} H^{n-1}(X)
= P^{n-1}(Y)/\operatorname{Im} P^{n-1}(X)
\simeq \operatorname{Ker}\bigl(\varphi_{*} \colon H^{n-1}(Y) \to H^{n+1}(X)\bigr),
\]
et l'on obtient la décomposition
\[
H^{*}(Y) \simeq \bigoplus_{\substack{i+j \leqslant n-1 \\ i,j \geqslant 0}}
L_{X}^{j}P^{i}(X) \ \oplus\ E^{n-1} .
\]
Sous cette identification, $\varphi_{*}$ agit comme $L_{X}$ sur le premier
morceau et par zéro sur $E^{n-1}$ ; $\Lambda_{Y}$ est nul sur $E^{n-1}$ et
coïncide avec $\Lambda_{X}$ sur l'autre. D'où le corollaire
\[
\operatorname{Ker}\varphi_{*} \simeq E^{n-1},
\qquad
\operatorname{Im}\varphi_{*} = \operatorname{Im}\varphi_{*}\varphi^{*}
= \operatorname{Im} L_{X} = \bigcap_{0 \leqslant i \leqslant n} \operatorname{Ker} p^{X}_{i} .
\]

\subsection*{La formule encadrée}

Le feuillet 39 conclut. En posant
$q = \mathrm{id}_{X} - \sum_{0 \leqslant i \leqslant 2n}p^{X}_{i}$, on a plus
généralement $\varphi_{*}\Lambda_{Y}^{m}\varphi^{*} = q\,\Lambda_{X}^{m-1}\,q$,
et dans l'autre sens
\[
\varphi^{*}\Lambda_{X}\varphi_{*} = \mathrm{id}_{Y} - p_{E},
\]
où $p_{E}$ est la projection sur le facteur $E^{n-1}$. Comme
$\Lambda_{Y}p_{E} = 0$, la multiplication à gauche par $\Lambda_{Y}$ élimine ce
terme, et il reste
\[
\boxed{\ \Lambda_{Y} = \varphi^{*}\,\Lambda_{X}^{2}\,\varphi_{*}\ }
\]
--- formule qui exprime le pseudo-inverse sur la section hyperplane à partir de
celui de l'ambiante. Elle donne aussitôt $C(X) \Rightarrow C(Y)$.

Il aborde alors la réciproque en introduisant les projections $p^{i,j}_{X}$ sur
les facteurs $L_{X}^{j}P^{i}(X)$, pour lesquelles
\[
\varphi_{*}\,p^{i,j}_{Y}\,\varphi^{*} = p^{i,j+1}_{X}
\qquad (0 \leqslant i+j \leqslant n-1),
\]
ce qui rend algébriques tous les $p^{i,j}_{X}$ avec $j \geqslant 1$, donc tous
sauf ceux de la première colonne. Le feuillet s'arrête sur un « Mais » : ce
sont précisément les $p^{i,0}_{X}$, en dimension médiane, qui résistent, et
c'est l'objet de la page suivante.

\pagerange{41}{41}
\section*{La réduction : $C(X)$ et l'algébricité de $\Lambda_{X}\varphi_{*}$ (page 41)}

Ce feuillet reprend au point où la page 39 s'arrête, et donne la réduction sous
sa forme la plus nette. \emph{Moyennant $C(Y)$}, la conjecture $C(X)$ équivaut
à chacun des énoncés suivants :

\begin{itemize}
  \item pour tout $i \leqslant n-1$, l'homomorphisme
  $H^{i}(Y) \to H^{i}(X)$, inverse à gauche de
  $\varphi^{i} \colon H^{i}(X) \to H^{i}(Y)$ et induit par
  $\Lambda_{X}\varphi_{*}$, est induit par une classe algébrique ;
  \item $\Lambda_{X}\varphi_{*}$ est algébrique.
\end{itemize}

L'opérateur $\Lambda_{X}\varphi_{*}$ correspond en effet à une classe de
cohomologie sur $X \times Y$, de degré $m - 1$ où $m = \dim(X \times Y) =
2n-1$, dans
\[
\sum_{i} H^{2(n-1)-i}(Y) \otimes H^{i}(X) = H^{2(n-1)}(X \times Y) .
\]
En lui appliquant l'opérateur de Lefschetz du produit,
$L_{X \times Y} = L_{X} \otimes 1 + 1 \otimes L_{Y}$, on obtient
\[
\Lambda_{X}\varphi_{*}L_{Y} + L_{X}\Lambda_{X}\varphi_{*}
= (\Lambda_{X}L_{X} + L_{X}\Lambda_{X})\,\varphi_{*}
= (\varphi_{*}\Lambda_{Y}\varphi^{*} + \mathrm{id}_{X})\,\varphi_{*},
\]
la dernière égalité étant la définition de $T_{X}$ du feuillet
32.\footnote{Une première rédaction du second membre, biffée, portait
$(2\,\mathrm{id} - \sum_{n \leqslant i \leqslant 2n}p^{X}_{i})\varphi_{*}$ ;
c'est la même chose au terme $p^{X}_{n}$ près, compté deux fois.} La classe de
$\Lambda_{X}\varphi_{*}$ est donc en dimension critique pour $X \times Y$, et
l'on obtient
\[
C(X) \iff \bigl[\,C(Y) \ \text{et}\ A^{0}(X \times Y)\,\bigr],
\]
puis, en itérant le long d'une suite décroissante de sections hyperplanes
$X \supset Y \supset Z \supset T \supset \cdots$,
\[
C(X) \implies A^{0}(X \times Y) \ \text{et}\ A^{0}(Y \times Z)
\ \text{et}\ A^{0}(Z \times T) \ \ldots
\]

\pagerange{43}{46}
\section*{La lettre du 4 janvier 1967 : où placer la conjecture de l'indice (pages 43 à 46)}

Le premier des deux carbones est une lettre de travail adressée à celui qui
rédigeait les notes d'un exposé sur les cycles algébriques.\footnote{L'en-tête
porte « Dear Coates » et la date « 4.1.67 » de sa main. Le second carbone porte
« 6.1.1966 » ; les deux dates sont transcrites telles quelles et rien sur les
feuillets ne les concilie. Les lettres grecques manquaient à la frappe et ont
été encrées à la main ; là où il ne les a pas encrées, l'indice reste absent.}

\subsection*{Ce qu'il demande qu'on écrive}

Sa recommandation est de rédaction avant d'être mathématique, et elle est
instructive : énoncer la conjecture de l'indice \emph{immédiatement après} les
résultats principaux de la théorie de Hodge, en avertissant qu'elle ne prendra
tout son sens qu'une fois couplée à la conjecture $A$ du paragraphe suivant.
Cela laisse ensuite la place d'exprimer les relations entre conjectures, du
type « $A$ + indice implique $B$ ».

Suit un inventaire de ce qui, en caractéristique zéro, se démontre : le
théorème de l'indice vaut ; les conditions $A_{\chi}$ à $D_{\chi}$ sont
indépendantes de la théorie cohomologique $\chi$, parce que l'existence de la
cohomologie de Betti les ramène toutes à leur analogue rationnel ; et $A$ et
$C$ sont indépendantes de la polarisation choisie --- $A$ parce qu'elle
équivaut à $B$, $C$ parce qu'elle s'exprime par les $A$ des produits successifs.
On peut donc, sans ambiguïté, écrire $A(X)$ à $D(X)$ sans indice.

En caractéristique $p > 0$ rien de tout cela n'est acquis, et il se corrige
explicitement : contrairement à ce qu'il avait affirmé trop vite pendant
l'exposé, sous l'influence de ses souvenirs du cas de caractéristique nulle, il
ne lui paraît pas clair que $A_{\chi}(X,\xi)$ et $C_{\chi}(X,\xi)$ soient alors
indépendantes de $\xi$ ; si la preuve ne se trouve pas, c'est la dépendance
possible qu'il faut signaler.\footnote{La lettre écrit ici « $C(X)$ is of
finite dimension over $\mathbb{Q}$ », où $C$ ne désigne plus la conjecture mais
le groupe des cycles modulo équivalence numérique ; et elle écrit
$C^{n-1}_{\chi}$ là où l'énoncé demande $C^{n-i}_{\chi}$ --- la condition $A$
s'exprimant par l'égalité $\dim C^{i} = \dim C^{n-i}$. Les deux glissements
sont sur le feuillet.}

\subsection*{La proposition}

Vient alors, en français dans le tapuscrit anglais, la proposition qui
rassemble les réductions des pages 32 à 41. Pour $X$ polarisée, avec $Y$ une
section hyperplane, les conditions suivantes sont équivalentes :

\begin{enumerate}
  \item[(i)] $C_{\chi}(X)$ ;
  \item[(ii)] $C_{\chi}(Y)$ et $A_{\chi}(X \times X)$ ;
  \item[(ii bis)] $C_{\chi}(Y)$ et $A_{\chi}(X \times X)^{0}$, l'exposant $0$
  signifiant qu'on se borne à la dimension critique
  $H^{2n-2} \to H^{2n+2}$ ;
  \item[(iii)] $C_{\chi}(Y)$ et $A_{\chi}(X \times Y)$ ;
  \item[(iii bis)] $C_{\chi}(Y)$ et $A_{\chi}(X \times Y)^{0}$, en dimension
  critique $H^{(2n-1)-1} \to H^{(2n-1)+1}$ ;
  \item[(iv)] $C_{\chi}(Y)$, et pour tout $i \leqslant n-1$ l'homomorphisme
  $H^{i}(Y) \to H^{i}(X)$ inverse à gauche de $\varphi^{i}$, induit par
  $\Lambda_{X}\varphi_{*}$, est induit par une correspondance algébrique ;
  \item[(iv bis)] $C_{\chi}(Y)$, et pour tout $j \geqslant n+1$ l'homomorphisme
  $H^{j}(X) \to H^{j-2}(Y)$ inverse à droite de $\varphi_{j-2}$, induit par
  $\varphi^{*}\Lambda_{X}$, est induit par une correspondance algébrique.
\end{enumerate}

Le \emph{corollaire} descend la chaîne : ces conditions équivalent aussi à
\[
A_{\chi}(X \times X)^{0} + A_{\chi}(Y \times Y)^{0} + A_{\chi}(Z \times Z)^{0} + \cdots
\]
pour une suite décroissante $X \supset Y \supset Z \supset \cdots$ de sections
hyperplanes, ou de même avec $X \times Y$, $Y \times Z$, etc. L'intérêt est
qu'on ramène ainsi la conjecture $A(k) = C(k)$ pour un corps entier à la seule
condition $A(T)^{0}$ pour toute variété $T$, c'est-à-dire à ce qui se passe en
dimension critique.

La démonstration tient en quelques observations. L'équivalence de (iv) et
(iv bis) est la transposition. Ces deux conditions entraînent (i) parce que
l'isomorphisme cherché $H^{2n-i}(X) \to H^{i}(X)$ se factorise

\begin{tikzcd}[column sep=small, nodes={font=\scriptsize}]
H^{2n-i}(X) \arrow[r, "\varphi^{*}\Lambda_{X}"] & H^{2n-i-2}(Y) \arrow[r, "\Lambda_{Y}^{(n-1)-i}"] & H^{i}(Y) \arrow[r, "\Lambda_{X}\varphi_{*}"] & H^{i}(X)
\end{tikzcd}

où les flèches extrêmes sont celles de (iv bis) et de (iv). Et elles sont
entraînées par (iii bis) grâce à la formule de la page 41,
\[
(\Lambda_{X}\varphi_{*})L_{Y} + L_{X}(\Lambda_{X}\varphi_{*})
= (\varphi_{*}\Lambda_{Y}\varphi^{*} + \mathrm{id}_{X})\,\varphi_{*},
\]
dont il note en marge la transposée
$(\varphi^{*}\Lambda_{X})L_{X} + L_{Y}(\varphi^{*}\Lambda_{X})
= \varphi^{*}(\varphi_{*}\Lambda_{Y}\varphi^{*} + \mathrm{id}_{X})$. Le reste
--- (iii) $\Rightarrow$ (iii bis), et (i) $\Rightarrow$ (iii) --- découle des
propriétés de stabilité.

\subsection*{L'état de la question, et une remarque d'histoire}

Il demande enfin qu'on dresse la liste de ce qui est su :

\begin{itemize}
  \item en caractéristique quelconque, $C(X)$ est connue si $\dim X \leqslant
  2$, et plus généralement $H^{2n-1}(X) \to H^{1}(X)$ est toujours induit par
  une classe de correspondance algébrique ; les projecteurs $\pi_{0}$,
  $\pi_{2n}$, $\pi_{1}$, $\pi_{2n-1}$ sont algébriques en toute dimension
  --- trivialement pour les deux premiers, moins trivialement pour les deux
  suivants ;
  \item en dimension $3$, on ignore, même en caractéristique nulle, si $C(X)$
  vaut ou seulement $D(X)$, et en caractéristique $p$ on ignore $A(X)$ et
  $B(X)$ ;
  \item en caractéristique $0$, $A(X)$ est connue pour $\dim X \leqslant 4$,
  mais $A(X)^{0}$ ne l'est pas en dimension $5$ ; le premier cas intéressant
  serait $X \times Y$ avec $\dim X = 3$ et $Y$ une section hyperplane, puisque
  cela donnerait $C(X)$.
\end{itemize}

« Ainsi les problèmes principaux se posent déjà pour les $1$-cycles sur les
variétés de dimension $3$, et partiellement même en caractéristique nulle. »

C'est ici qu'il glisse la remarque qui vaut d'être citée. L'algébricité des
$\pi_{i}$ pour une surface, lui apprend Tate, avait été signalée par Hodge dans
un article de 1938 ; et il ajoute qu'il est « assez frappant que cet énoncé
n'ait pas davantage frappé les géomètres algébristes, et qu'il soit tombé dans
l'oubli pendant près de trente ans ! ».\footnote{La référence est donnée par la
lettre elle-même et n'a pas été vérifiée ici : Hodge, \emph{Algebraic
correspondences between surfaces}, Proc.\ London Math.\ Soc., série 2, vol.\
XLIV, 1938, p.\ 226.}

La lettre se termine sur un projet et sur son abandon. Poussé par une question
de Kleiman, il annonce qu'il va reprendre ses « vieux gribouillages » sur la
réduction de la forme forte de Lefschetz à la forme faible --- ce sont les
feuillets 50 à 54. Et il enterre la suggestion qu'il avait faite au téléphone,
de ramener toute variété à une hypersurface par équivalence birationnelle et
spécialisation : Serre lui a fait observer qu'une telle variété serait
nécessairement simplement connexe. Il faudra donc travailler aussi avec des
hypersurfaces singulières, « et voir comment reformuler les conjectures
standard pour les variétés singulières » --- ce que le dossier ne fait pas.

\pagerange{47}{49}
\section*{La lettre du 6 janvier : Cayley--Hamilton, et les variantes « faibles » (pages 47 à 49)}

Le second carbone fait deux choses : il rend $C_{\chi}(X)$ indépendante de la
polarisation par un argument purement linéaire, et il montre que des variantes
apparemment plus faibles des conjectures les impliquent en fait toutes.

\subsection*{La proposition, et sa preuve}

\emph{$C(X)$ équivaut à chacune des conditions suivantes :}

\begin{enumerate}
  \item[a)] $D_{\chi}(X)$ vaut, et pour tout $i < n$ il existe un isomorphisme
  algébrique $H^{2n-i}(X) \to H^{i}(X)$ --- sans qu'on affirme rien de ce qu'il
  induit dans les autres degrés ;
  \item[b)] pour tout endomorphisme algébrique de $H^{i}(X)$ les coefficients
  du polynôme caractéristique sont rationnels, et pour tout $i < n$ il existe
  un isomorphisme algébrique $H^{2n-i}(X) \to H^{i}(X)$.
\end{enumerate}

L'intérêt de l'énoncé est qu'il ne demande qu'un isomorphisme algébrique
\emph{quelconque}, et non l'inverse de $L^{n-i}$ ; il montre, dit-il, le rôle
joué par le fait que $H^{i}(X)$ doive être « motiviquement » isomorphe à son
dual naturel $H^{2n-i}(X)$.

La démonstration de b) $\Rightarrow$ $C_{\chi}(X)$ est l'argument le plus court
du dossier. Soit $u \colon H^{2n-i}(X) \to H^{i}(X)$ l'isomorphisme algébrique
donné, et $v \colon H^{i}(X) \to H^{2n-i}(X)$ celui induit par $L_{X}^{n-i}$,
algébrique lui aussi. Alors $w = uv$ est un automorphisme algébrique de
$H^{i}(X)$, et la relation de Cayley--Hamilton
\[
w^{b} - \sigma_{1}(w)\,w^{b-1} + \cdots + (-1)^{b}\sigma_{b}(w) = 0,
\qquad b = \operatorname{rang} H^{i}(X),
\]
exprime $w^{-1}$ comme combinaison linéaire des $w^{i}$ à coefficients
$\pm\,\sigma_{i}(w)/\sigma_{b}(w)$. L'hypothèse rend ces coefficients
rationnels, donc $w^{-1}$ est algébrique, donc $w^{-1}u = v^{-1}$ aussi ---
c'est-à-dire l'inverse de $L_{X}^{n-i}$.\footnote{Le tapuscrit écrit la
relation de Cayley--Hamilton en $u$ alors que ses coefficients sont ceux de
$w$ ; elle est rétablie en $w$ ici. En caractéristique nulle, note-t-il,
l'énoncé se simplifie : $C(X)$ équivaut à la seule existence d'isomorphismes
algébriques $H^{2n-i}(X) \to H^{i}(X)$, la condition préliminaire de b) étant
alors automatique.}

Le \emph{corollaire} est immédiat : si $X$ et $X'$ satisfont $C_{\chi}$ et si
$u \colon H^{i}(X) \to H^{i+2j}(X')$ est un isomorphisme algébrique, alors
$u^{-1}$ est algébrique --- car les deux espaces s'identifient algébriquement
à leurs duaux dans les deux sens, si bien que la transposée de $u$ se relit
$u' \colon H^{i+2j}(X') \to H^{i}(X)$, et $u'u$ est un automorphisme
algébrique auquel on applique ce qui précède.

C'est ici qu'il formule le programme dans toute son ampleur, et aussitôt le
ramène à terre. Le même résultat devrait valoir pour $u$ seulement
\emph{mono}morphique, à condition de disposer d'un inverse à gauche algébrique
--- ce qu'on a bien pour $H^{n-1}(X) \to H^{n-1}(Y)$, où l'inverse est
$\Lambda_{X}\varphi_{*}$ --- mais il semble qu'il faille en général la relation
d'indice de Hodge en plus. « Le yoga complet étant alors qu'on a la catégorie
des motifs, qui est semi-simple ! » Puis : « sans parler de motifs, et en
restant les pieds sur terre », il serait bon d'expliquer dans les notes que
$C(X)$ jointe à la relation d'indice $I(X \times X)$ entraîne que l'anneau des
classes de correspondances de $X$ est semi-simple, et d'en déduire l'existence
des inverses à gauche et à droite cherchés.\footnote{La catégorie des motifs
purs n'est ici qu'un mot d'ordre : rien sur ces feuillets n'en construit ni
n'en suppose la définition, et la semi-simplicité y est présentée comme une
conséquence attendue, non comme un acquis.}

\subsection*{Les variantes faibles, qui ne le sont pas}

La lettre finit sur une remarque qu'il annonce triviale et qui ne l'est pas.
Introduisons
\begin{itemize}
  \item $A'_{\chi}(X)$ : si $2i \leqslant n-1$, tout élément de $H^{i}(X)$ dont
  l'image dans $H^{i}(Y)$ est algébrique est algébrique ;
  \item $A''_{\chi}(X)$ : si $2i \geqslant n-1$, tout élément algébrique de
  $H^{i+2}(X)$ est l'image d'un élément algébrique de $H^{i}(Y)$ ;
\end{itemize}
et leurs versions $A'^{0}_{\chi}$, $A''^{0}_{\chi}$ restreintes à la dimension
critique ($2i = n-1$ si $n$ est impair, $2i = n-2$ s'il est pair). Ce sont des
énoncés « à la Lefschetz faible », trivialement impliqués par $A_{\chi}(X)$
et par $C_{\chi}(Y)$ respectivement --- dans le premier cas en appliquant
$\varphi$, dans le second en prenant $y = \Lambda_{Y}\varphi^{*}(x)$.

Or ces variantes prétendument faibles impliquent les relations de Lefschetz
complètes :
\[
C_{\chi}(X) \iff
C_{\chi}(Y) + A'^{0}_{\chi}(X \times X) + A''^{0}_{\chi}(X \times X),
\]
et donc, par récurrence, à la conjonction des $A'^{0}_{\chi}$ et
$A''^{0}_{\chi}$ pour toutes les variétés $X \times X$, $Y \times Y$,
$Z \times Z$, \ldots (ou $X \times Y$, $Y \times Z$, \ldots).

Tout vient de ce que $A_{\chi}(T)^{0}$ résulte de la conjonction de
$A'_{\chi}(T)^{0}$ et $A''_{\chi}(T)^{0}$, comme on le voit en décomposant
l'opérateur de la dimension critique à travers la section hyperplane. Si
$\dim T = 2m$ est pair, $L_{T}^{2} \colon H^{2m-2}(T) \to H^{2m+2}(T)$ se
factorise

\begin{tikzcd}[column sep=small, nodes={font=\scriptsize}]
H^{2m-2}(T) \arrow[r, "\varphi^{*}"] & H^{2m-2}(Y) \arrow[r, "\varphi_{*}"] & H^{2m}(T) \arrow[r, "L_{T}"] & H^{2m+2}(T)
\end{tikzcd}

et si $\dim T = 2m+1$ est impair, $H^{2m}(T) \to H^{2m+2}(T)$ se factorise

\begin{tikzcd}[column sep=small, nodes={font=\scriptsize}]
H^{2m}(T) \arrow[r, "\varphi^{*}"] & H^{2m}(Y) \arrow[r, "\varphi_{*}"] & H^{2m+2}(T)
\end{tikzcd}

\pagerange{50}{54}
\section*{Le lemme : un anneau d'endomorphismes qui contient $\Lambda$ (pages 50, 52 et 54)}

Les trois derniers feuillets sont les « vieux gribouillages » que la lettre
annonçait : la réduction de la forme forte de Lefschetz à la forme faible,
menée cette fois au niveau de l'anneau. C'est le théorème du n° 6 transporté
dans l'anneau des correspondances algébriques.

\emph{Lemme.} Soit $A$ un anneau d'endomorphismes du groupe abélien
$H^{*}(X)$, stable par composantes homogènes et par transposition. Supposons
$\zeta \in A$, et supposons que pour tout $0 \leqslant i < n$ il existe
$u_{i} \in A$ tel que le composé
\[
H^{i}(X) \xrightarrow{\ \zeta^{n-i}\ } H^{2n-i}(X) \xrightarrow{\ u_{i}\ } H^{i}(X)
\]
soit l'identité. Alors $A$ contient le pseudo-inverse $\Lambda$ de $\zeta$ et
les projections sur les facteurs élémentaires
$\zeta^{j}P^{i}$.\footnote{Le feuillet écrit « Alors $A$ contient $L$ », la
lettre étant tracée comme l'opérateur de Lefschetz du premier cahier. Comme
$\zeta$ est dans $A$ par hypothèse, et que $\zeta$ est l'opérateur de Lefschetz
de l'énoncé, la seule conclusion qui apporte quelque chose est celle portant
sur le pseudo-inverse : c'est ainsi qu'on l'énonce ici. Le feuillet ne tranche
pas, et l'apparat de la transcription conserve sa lettre.} L'hypothèse de
stabilité par composantes homogènes tient exactement le rôle que jouait, au
n° 5, le sur-anneau $\mathcal{F}$ où les projecteurs vivaient d'avance : elle
donne les projections sur les $H^{i}$ sans qu'on ait à les supposer.

\emph{Corollaire.} Sont équivalentes : (i) $\Lambda$ est algébrique ; (ii) pour
$1 \leqslant i \leqslant n$, il existe une correspondance algébrique induisant
l'inverse $H^{2n-i}(X) \to H^{i}(X)$ de $\zeta^{n-i}$. Et s'il en est ainsi,
les projections sur les facteurs élémentaires de $H^{*}(X)$ sont algébriques,
ainsi que les projections $H^{*}(X) \to H^{i}(X)$ --- c'est-à-dire que la
conjecture de type Lefschetz entraîne l'algébricité des projecteurs de
Künneth.

\subsection*{La démonstration, et sa lacune}

Par stabilité sous transposition, il suffit de traiter les $\zeta^{j}P^{i}$
avec $i + 2j \leqslant n$. On pose $k = i + 2j$ et l'on raisonne par récurrence
sur $k$, le cas $k < 0$ étant vide. Supposons donc le résultat acquis pour
$k' < k$ : on sait déjà que les projections $p_{i}$ sur $H^{i}$ sont dans $A$
pour $i < k$, donc par transposition pour $i > 2n-k$. On introduit
\[
T_{k} = \mathrm{id} - \sum_{i < k} p_{i} - \sum_{j > 2n-k} p_{j} \ \in A,
\]
la projection sur la bande des degrés $k \leqslant \alpha \leqslant 2n-k$, et
l'on obtient les projections cherchées par
\[
p_{k} = u_{k}\,\zeta^{n-k}\,T_{k},
\qquad
\text{puis} \qquad
\zeta^{j}\,u_{k-2j}\,\zeta^{n-k+j}\,T_{k}
\]
pour la projection sur $\zeta^{j}H^{k-2j}$, dont on déduit par différence,
grâce à une récurrence ascendante sur $i$ (descendante sur $j$), la projection
sur chaque $\zeta^{j}P^{i}$ avec $i + 2j = k$.

Il faut dire où cette démonstration, telle que les feuillets la portent, ne se
referme pas. Le pas décisif est l'affirmation que $\zeta^{n-k}T_{k}$ est nul
sur $H^{\alpha}$ pour $\alpha \neq k$ et induit $\zeta^{n-k}$ sur $H^{k}$ ---
ce qui donnerait $u_{k}\zeta^{n-k}T_{k} = p_{k}$. Or $T_{k}$ retient toute la
bande $k \leqslant \alpha \leqslant 2n-k$, et $\zeta^{n-k}$ n'y est pas nul :
il envoie $H^{\alpha}$ dans $H^{\alpha+2n-2k}$, qui est non nul dès que
$\alpha \leqslant 2k$. L'annulation demandée n'est donc pas acquise par la
seule définition de $T_{k}$, et rien sur les trois feuillets ne la
justifie.\footnote{Ces pages sont écrites très vite : les formules portent, la
prose qui les relie souvent non, et une insertion interligne fait intervenir un
objet $N$ qui n'est défini nulle part. Il se peut donc que la justification ait
été sur le feuillet et se soit perdue à la lecture ; on enregistre la lacune
telle qu'elle se présente, sans la combler et sans la mettre au compte de
l'argument.}

Le dossier s'arrête là, au tiers d'un feuillet, sur les mots « Ceci prouve le
lemme ».

\end{document}
